% Every operation sgt offers, grouped by what it does to the recorded set, with
% the state before and after each command and what each one costs.
% Prints as Figure 5, after the fork figure in Section 4.3.
\newcommand{\FigFourCaption}{The nine operations that act on the record, grouped
by what each one does to it. One command writes records and the other eight read
them or change which of them a version includes, and no argument anywhere in the
figure is a hash or a line number: every one of them is a request the developer
named, a branch, or a function.}

\newcommand{\FigFourDescription}{Three panels, one per group of operations. The
first panel holds the single operation that writes records, the second holds the
three that only read them, and the third holds the five that change which
records a version includes. Each row carries a small diagram of the set before
and after the command, the command itself, one sentence saying what the developer
gets, and one sentence in a grey band on the right saying what the operation
costs. Each diagram draws one horizontal line per request with a mark for every
edit filed under that request, as in Figure 1, and a dotted outline is an edit
the version no longer includes.}

\begin{tikzpicture}[x=1mm, y=1mm, line cap=round, font=\sffamily]
% narrow ragged columns, so break a line early rather than hyphenate a word
\hyphenpenalty=10000 \exhyphenpenalty=10000

% the thumbnails carry the before/after convention, so the column says so itself
\node[anchor=west, inner sep=0, text=sgtInk!60,
      font=\sffamily\fontsize{5.0}{5.8}\selectfont] at (3.4,104.6)
  {before \ttarrow{} after};
\node[hd, anchor=west] at (19,104.6)    {COMMAND};
\node[hd, anchor=west] at (53,104.6)    {WHAT THE DEVELOPER GETS};
\node[hd, anchor=west] at (115.5,104.6) {WHAT IT COSTS};

% one group of operations. The cost column gets a band of its own on the right,
% the same band Figure 2 uses, so the reader can read that column straight down.
% #1 y, #2 height, #3 tab colour, #4 tab text
\newcommand{\opgroup}[4]{%
  \panel{0}{#1}{176}{#2}{#3}{#4}
  \fill[wGrey, rounded corners=0.8mm] (112.5,{#1+1.2}) rectangle (174.6,{#1+#2-1.2});}

% #1 y of the first text line, #2 thumbnail, #3 command, #4 gets, #5 costs
\newcommand{\oprow}[5]{%
  \node[inner sep=0] at (10,{#1-2.1}) {#2};
  \node[anchor=north west, text width=32mm, align=left, inner sep=0, text=sgtInk,
        font=\ttfamily\fontsize{5.6}{6.6}\selectfont] at (19,#1) {#3};
  \node[anchor=north west, text width=55mm, align=left, inner sep=0, text=sgtInk!85,
        font=\sffamily\fontsize{5.7}{6.7}\selectfont] at (53,#1) {#4};
  \node[anchor=north west, text width=56mm, align=left, inner sep=0, text=sgtInk!85,
        font=\sffamily\fontsize{5.7}{6.7}\selectfont] at (115.5,#1) {#5};}

% =================== what writes records ==================================
% One row, and the group holding one row is the point: the developer types one
% command to add to the record, and everything else in the figure decides what
% that record is read for or what a version does with it.
\opgroup{86.4}{13.6}{sgtInk}{WHAT WRITES RECORDS}

% Centred in its panel rather than hung from the top of it, so that the space
% around the one row reads as emphasis and not as a row somebody took out.
\oprow{94.3}{\thsave}{sgt save -m "..."}
  {One record per changed function, filed under a name for the request. The only
   command that adds to the record.}
  {The grouping is a guess. Correcting it moves labels and changes no recorded
   edit.}

% =================== what only reads them =================================
\opgroup{55.4}{29}{sgtGrey}{WHAT ONLY READS THEM}

\oprow{79.8}{\thwhy}{sgt show api.py::handle}
  {The request this function came from, the saves that touched it, and how much
   other work removing it takes along.}
  {The answer is exactly as good as the grouping behind it.}
\oprow{71.4}{\thdiff}{sgt diff main feat/export}
  {The edits each version has that the other does not, in both directions.}
  {It answers in functions, so seeing which lines moved still means opening the
   file.}
\oprow{63.0}{\thckpt}{sgt intent list}
  {Each feature cut into the stretches of work done on it, each with a name like
   {\ttfamily rate limiting@2} the other commands accept.}
  {A language model places the cuts by reading the saves, so a stretch can start
   one save too early.}

% ============ what changes which records a version includes ===============
\opgroup{7.4}{46}{sgtInk}{WHAT CHANGES WHICH RECORDS A VERSION INCLUDES}

\oprow{48.8}{\threvert}{sgt revert "price cache"}
  {The named set subtracted from the version, with every function it changes
   named before anything is applied.}
  {A function later work also edited gets the set's lines subtracted from its
   current text, and the tests decide.}
\oprow{40.4}{\threstore}{sgt restore "price cache"}
  {The same set back, byte for byte, with everything added since left alone.}
  {Refused when putting the set back would leave two versions of one function
   live, and the refusal prints edit identifiers with no function name.}
\oprow{32.0}{\thswitch}{sgt switch feat/export}
  {The working files rebuilt from that branch's set.}
  {Refused on unsaved edits, because \sgt{} has nowhere to put them yet.}
\oprow{23.6}{\thfork}{sgt resolve api.py::handle}
  {The two competing versions side by side, and whichever one the developer
   picks or writes is included once the project's test command passes.}
  {Two edits at opposite ends of one function are also competing versions, and
   git would have merged them without asking.}
\oprow{15.2}{\thsubset}{sgt propose land --subset "export"}
  {One finished feature sent to the shared branch, with the unfinished one left
   behind.}
  {Two functions that sit next to each other in one file can lose the blank line
   between them.}

\end{tikzpicture}
